%% ---------------------------------------------------------------------------
%% The linear-logic tensors: \ltimes and \rtimes as Dialectica products
%% A section of "The Category of Containers" (MacBeth)
%% MacBeth (Kodamai AI Mathematician project), 2026-07-17
%%
%% Written to be DROPPED INTO the four-monoidal chapter:
%%   * delete the preamble down to \begin{document} (the chapter already has the
%%     macros \Set \Cont \Poly \Fam \Ext \dir \tri \prov and the
%%     Theorem/Definition/... environments);
%%   * the top-level \section{...} below becomes a chapter section.
%% Compiles standalone with pdflatex as it is (run twice for references).
%% ---------------------------------------------------------------------------
\documentclass[11pt,a4paper]{article}

\usepackage[margin=2.6cm]{geometry}
\usepackage{amsmath,amssymb,amsthm}
\usepackage{stmaryrd}   %% \llbracket, \rrbracket for the extension brackets
\usepackage{booktabs}

%% ---- Notation (kept consistent with the four-monoidal chapter) ----
\newcommand{\Set}{\mathbf{Set}}
\newcommand{\Cont}{\mathbf{Cont}}
\newcommand{\Poly}{\mathbf{Poly}}
\newcommand{\Fam}{\mathbf{Fam}}
\newcommand{\Ext}[1]{\llbracket #1 \rrbracket}
\newcommand{\id}{\mathrm{id}}
\newcommand{\dir}{\mathbin{\otimes}}          %% Dirichlet tensor
\newcommand{\tri}{\mathbin{\triangleleft}}    %% sequential operator
\newcommand{\Dial}{\mathbf{Dial}}
\newcommand{\Hmg}{\mathrm{Hmg}}
\newcommand{\dtens}{\mathbin{\otimes_{\Dial}}} %% de Paiva's Dialectica tensor

%% ---- Flagging environments (Neil's convention) ----
\theoremstyle{plain}
\newtheorem{Theorem}{Theorem}[section]
\newtheorem{Proposition}[Theorem]{Proposition}
\newtheorem{Lemma}[Theorem]{Lemma}
\newtheorem{Corollary}[Theorem]{Corollary}
\newtheorem{Conjecture}[Theorem]{Conjecture}
\theoremstyle{definition}
\newtheorem{Definition}[Theorem]{Definition}
\newtheorem{Example}[Theorem]{Example}
\theoremstyle{remark}
\newtheorem{Remark}[Theorem]{Remark}
\newtheorem{Question}[Theorem]{Question}

\newcommand{\prov}[1]{\hfill\textsf{\footnotesize[#1]}}

\title{\textbf{The linear-logic tensors: $\ltimes$ and $\rtimes$ as
  Dialectica products}}
\author{MacBeth\\ \normalsize Kodamai AI Mathematician project}
\date{17 July 2026}

\begin{document}
\maketitle

%% ===========================================================================
\section{The linear-logic tensors: $\ltimes$ and $\rtimes$ as Dialectica
  products}
\label{sec:dialectica}

\subsection{Setup: is there anything outside the Day family?}
\label{sec:dia-setup}

Recall the standing conventions. $\Cont$ is the category of containers,
equivalent to the polynomial functors $\Poly$; a container $p=(S_p,p[-])$ has a
set $S_p$ of shapes and, at each shape $s$, a set $p[s]$ of directions
(positions), with extension $\Ext pX=\sum_{s\in S_p}X^{p[s]}$. We write
$y^A$ for the representable with one shape and directions $A$, and $y:=y^1$ for
the unit container. Four monoidal structures are the standard census: the
coproduct $+$, the categorical product $\times$, the Dirichlet tensor $\dir$,
and the sequential (substitution) operator $\tri$.

The classification theorem (Theorem~A) accounts for the \emph{Day-convolutional}
part of this census. Call a monoidal structure $(\odot,J)$ on $\Cont$
\emph{convolutional} when $\odot$ preserves coproducts in each variable and
carries representables to representables; equivalently
$(p\odot q)[(s,t)]=p[s]\star q[t]$ for a monoidal structure $(\star,I)$ on
$\Set$, of which $\odot$ is the Day convolution. Theorem~A says the Day
construction $(\star,I)\mapsto(\odot_\star,y^I)$ is an \emph{equivalence} between
monoidal structures on $\Set$ and convolutional monoidal structures on $\Cont$;
$\times$ and $\dir$ are the images of $(+,\varnothing)$ and $(\times,1)$, while
$+$ and $\tri$ sit outside for elementary reasons.

An equivalence onto the \emph{convolutional} structures is, honestly read, a
statement about a sub-family. It leaves a question open at the level of $\Cont$
itself:

\begin{quote}
\emph{Are there monoidal structures on $\Cont$ outside the Day family
altogether --- tensors whose direction at $(s,t)$ is \textbf{not} a pointwise
function $F(p[s],q[t])$ of the two fibres?}
\end{quote}

The answer is yes, and the examples are not exotic: they were written down by
Dorta--Jarvis--Niu \cite{djn2305} and left \emph{uninterpreted}. This section
identifies them as the Dialectica tensors of Gödel's functional interpretation,
answering DJN's own stated open question.

\subsection{The two DJN tensors in container coordinates}
\label{sec:dia-def}

Dorta--Jarvis--Niu study monoidal structures on generalised polynomial
categories $\Sigma\Pi\mathcal C$; specialising to $\mathcal C=\mathbf 1$ gives
$\Sigma\Pi\mathbf 1\simeq\Poly\simeq\Cont$ \cite[Prop.~2.9]{djn2305}. Beyond the
Day convolution $\dir$ and the composition product $\tri$, their \S6 exhibits
\emph{at least two further} monoidal structures with the same unit $y$, which we
write $\ltimes$ and $\rtimes$.

\begin{Definition}[$\ltimes$ and $\rtimes$, in container coordinates]
\label{def:dia-tensors}
Both tensors have shape set $S_p\times S_q$. Their directions are
\[
  (p\ltimes q)[(s,t)] \;=\; p[s]^{\,S_q}\times q[t]^{\,S_p},
  \qquad
  (p\rtimes q)[(s,t)] \;=\; p[s]^{\,S_q}\times q[t],
\]
with unit $y$ on both sides. Here $X^{S}$ denotes the function set $S\to X$, and
the exponent $S$ is the \emph{shape set of the other factor}.
\prov{Cited: DJN \cite[\S6]{djn2305}, in container coordinates}
\end{Definition}

The tensor $\ltimes$ is \emph{symmetric}: each factor's direction set is
exponentiated by the other factor's shape set. The tensor $\rtimes$ is
\emph{triangular}: only the left factor is exponentiated, so it will turn out to
be associative but not symmetric (Lemma~\ref{lem:nfold}). These are two of the
monoidal products DJN define; they explicitly leave their meaning open, writing
that ``we would like to know if there are interpretations or applications for
these monoidal products'' \cite{djn2305}.

The single feature that governs everything below is the exponential
$(-)^{S_q}$: a direction set exponentiated by a \emph{global} shape set, rather
than the pointwise $p[s]\star q[t]$ of a Day tensor. That is what pushes
$\ltimes/\rtimes$ outside the Day family (§\ref{sec:dia-consequences}), and it
is the Dialectica hallmark (§\ref{sec:dia-main}).

\begin{Lemma}[Unit laws and closed forms]
\label{lem:nfold}
$p\ltimes y\cong p\cong y\ltimes p$ and $p\rtimes y\cong p\cong y\rtimes p$.
The $k$-fold iterates have closed forms
\begin{align*}
  (p_1\ltimes\dots\ltimes p_k)\bigl[(s_1,\dots,s_k)\bigr]
    &\;\cong\; \prod_{f=1}^{k} p_f[s_f]^{\,\bigl(\prod_{g\neq f}S_{p_g}\bigr)},\\
  (p_1\rtimes\dots\rtimes p_k)\bigl[(s_1,\dots,s_k)\bigr]
    &\;\cong\; \prod_{f=1}^{k} p_f[s_f]^{\,\bigl(\prod_{g> f}S_{p_g}\bigr)} .
\end{align*}
Consequently $\ltimes$ is associative and symmetric, while $\rtimes$ is
associative but not symmetric.
\prov{MacBeth (computed; exponential laws)}
\end{Lemma}

\begin{proof}
For the unit, $S_y=1$ and $y[\ast]=1$, so
$(p\ltimes y)[(s,\ast)]=p[s]^{1}\times 1^{S_p}=p[s]$, and symmetrically; the
$\rtimes$ case is the same computation. The closed forms follow by induction from
Definition~\ref{def:dia-tensors} using $(X\times Y)^{S}\cong X^{S}\times Y^{S}$
and $X^{S\times T}\cong(X^{S})^{T}$. In the $\ltimes$ form each factor $p_f$ is
exponentiated by the product of \emph{all} other shape sets, an expression
invariant under reassociation and under permutation of the factors; hence
$\ltimes$ is associative and symmetric. In the $\rtimes$ form each factor is
exponentiated only by the shape sets \emph{to its right}, an expression
unambiguous under reassociation (the product over $g>f$ does not depend on
bracketing) but sensitive to the factor order; hence $\rtimes$ is associative but
not symmetric.
\end{proof}

The triangular exponent pattern is the precise sense in which $\rtimes$ is
\emph{directed}: reading the factors left to right, factor $f$ is challenged by
the choices of all \emph{later} factors, but not the earlier ones --- a strict,
one-directional dependency, of which $\ltimes$ is the symmetric closure.

\subsection{The homogeneous fragment is de Paiva's Dialectica category}
\label{sec:dia-hmg}

\begin{Definition}[Homogeneous container]
\label{def:homogeneous}
A container is \emph{homogeneous} of arity $(I,A)$ if it has the form
$p=\sum_{i\in I}y^{A}$: shape set $I$, with \emph{every} direction set equal to
the same $A$.
\prov{Cited: DJN \cite[\S2]{djn2305}}
\end{Definition}

DJN prove that the homogeneous objects form a full subcategory equivalent to de
Paiva's original Dialectica category
\cite[Prop.~2.13]{djn2305}:
\[
  \Hmg(2)\;\simeq\;\Dial(\Set).
\]
Under this dictionary a homogeneous object corresponds to a de Paiva triple
$(U,X,\alpha)$, where $U$ is the set of \emph{positions} (the shape set $I$),
$X$ is the set of \emph{challenges} (the common direction set $A$), and $\alpha$
is a set-valued (relational) predicate on $U\times X$. A morphism carries a
forward map $f\colon U\to V$ together with a backward map $F\colon U\times Y\to
X$ --- exactly the two components of a container morphism between homogeneous
objects. So \emph{positions are $U$ and directions are $X$}.

\begin{Definition}[de Paiva's Dialectica tensor]
\label{def:depaiva-tensor}
On $\Dial(\Set)$ the multiplicative (linear-logic) tensor is
\[
  (U,X,\alpha)\dtens(V,Y,\beta)
    \;=\;\bigl(\,U\times V,\ X^{V}\times Y^{U},\ \alpha\otimes\beta\,\bigr),
\]
where, for $f\colon V\to X$ and $g\colon U\to Y$, the predicate is
$(\alpha\otimes\beta)\bigl((u,v),(f,g)\bigr)=\alpha(u,f v)\wedge\beta(v,g u)$.
\prov{Cited: de Paiva \cite{depaiva89}}
\end{Definition}

The witness object $X^{V}\times Y^{U}$ is the defining feature of the
Gödel--de~Paiva functional interpretation: to answer a challenge against a
conjunction $A\otimes B$ one must supply, for each of the opponent's moves in
one component, a response in the \emph{other} --- hence the mutual function
spaces.

\subsection{Main identification: $\ltimes$ is de Paiva's tensor, extended}
\label{sec:dia-main}

\begin{Proposition}[$\ltimes$ restricts to the Dialectica tensor]
\label{prop:ltimes-is-dial}
On the homogeneous fragment $\Hmg(2)\simeq\Dial(\Set)$, the tensor $\ltimes$
restricts to de Paiva's Dialectica tensor $\dtens$. Concretely, for homogeneous
$p=\sum_{I}y^{A}$ and $q=\sum_{J}y^{B}$,
\[
  p\ltimes q \;=\;\sum_{I\times J} y^{\,A^{J}\times B^{I}},
\]
whose shape set $I\times J=U\times V$ and direction set
$A^{J}\times B^{I}=X^{V}\times Y^{U}$ match Definition~\ref{def:depaiva-tensor}
exactly. Hence $\ltimes$ is de Paiva's Dialectica tensor extended off
$\Hmg(2)$ to all of $\Poly\simeq\Cont$.
\prov{MacBeth (an identification; matches definitions)}
\end{Proposition}

\begin{proof}
Take $p$ and $q$ homogeneous with $S_p=I$, $p[i]=A$ for all $i$, and $S_q=J$,
$q[j]=B$ for all $j$. Definition~\ref{def:dia-tensors} gives shape set
$I\times J$ and, at every shape $(i,j)$,
\[
  (p\ltimes q)[(i,j)]\;=\;p[i]^{\,S_q}\times q[j]^{\,S_p}
    \;=\;A^{J}\times B^{I}.
\]
This direction set is independent of $(i,j)$, so $p\ltimes q$ is again
homogeneous, of arity $\bigl(I\times J,\ A^{J}\times B^{I}\bigr)$, i.e.\
$p\ltimes q=\sum_{I\times J}y^{A^{J}\times B^{I}}$. Under the dictionary of
§\ref{sec:dia-hmg} --- $U=I$, $V=J$, $X=A$, $Y=B$ --- the shape set is
$U\times V$ and the direction set is $X^{V}\times Y^{U}$, the forward and
backward objects of $\dtens$ in Definition~\ref{def:depaiva-tensor}. The
predicate matches by the same substitution: a direction $(f,g)\in A^{J}\times
B^{I}$ is precisely a pair $f\colon V\to X$, $g\colon U\to Y$, and DJN's fibrewise
predicate composition realises $\alpha\otimes\beta$. So on $\Hmg(2)$ the two
tensors agree, and $\ltimes$ is the extension of $\dtens$ to all containers.
\end{proof}

For $\rtimes$ the corresponding restriction gives, on homogeneous inputs,
direction $A^{J}\times B$: the left challenge is functionalised over $V$ but the
right challenge is presented bare. By Lemma~\ref{lem:nfold} the $k$-fold
$\rtimes$ exponentiates each factor's challenge only over the positions of the
factors to its \emph{right}. So $\rtimes$ is the \emph{directed} Dialectica
tensor: where $\ltimes$ is the conjunction in which both players respond
adaptively to the other's move, $\rtimes$ is the conjunction whose dependency
flows one way only --- one problem answered adaptively in response to the other,
which is played blind. Contrast the two closed forms of Lemma~\ref{lem:nfold}:
$\ltimes$ exponentiates over the product of \emph{all} other shape sets and is
symmetric; $\rtimes$ exponentiates only rightward and is directed.

For orientation, here is how the two Dialectica conjunctions split between the
Cont tensors. On $\Hmg(2)$ the Day tensor $\dir$ restricts to direction
$A\times B=X\times Y$ --- the naive Gödel conjunction $A\wedge B$, whose
challenge is the plain product of challenges --- whereas $\ltimes$ restricts to
the linear tensor $X^{V}\times Y^{U}$. The two structures cleanly separate the
non-linear ($\dir$) and multiplicative-linear ($\ltimes$) conjunctions of the
Dialectica interpretation.

\subsection{Structural consequences}
\label{sec:dia-consequences}

\begin{Proposition}[Non-convolutional]
\label{prop:nonconv}
Neither $\ltimes$ nor $\rtimes$ is convolutional: there is no functor
$F\colon\Set\times\Set\to\Set$ with $(p\ltimes q)[(s,t)]=F(p[s],q[t])$. In
particular they are not Day convolutions of any monoidal structure on $\Set$, so
Theorem~A does not reach them.
\prov{MacBeth (computed witness)}
\end{Proposition}

\begin{proof}
Let $p=y^{2}$ (one shape, direction set of size $2$) and, for each $n$, let
$q_n$ be homogeneous with $|S_{q_n}|=n$ and every direction set of size $2$.
Then $(p\ltimes q_n)[(\ast,t)]=2^{\,n}\times 2^{\,1}=2^{\,n+1}$. The two fibres
$p[\ast]=2$ and $q_n[t]=2$ are fixed independently of $n$, yet the resulting
direction set grows with $n=|S_{q_n}|$. A convolutional formula $F(2,2)$ would be
constant in $n$; no such $F$ exists. The same inputs defeat $\rtimes$. Since
every Day tensor is convolutional (with $F=\star$), neither tensor is a Day
convolution, and Theorem~A --- an equivalence onto the \emph{convolutional}
structures --- does not apply to them.
\end{proof}

This is the honest boundary of the classification. The Day family --- the four
canonical tensors together with the wider family indexed by monoidal structures
on $\Set$ --- does \emph{not} exhaust the monoidal structures on $\Cont$. The
structures it omits are the linear-logic tensors $\ltimes$ and $\rtimes$, and
the reason it omits them is precisely that their direction is a non-pointwise,
shape-set-indexed exponential rather than a fibrewise formula.

\begin{Proposition}[Non-cocontinuous]
\label{prop:noncocont}
$(-)\ltimes q$ does not preserve coproducts: $\ltimes$ does not distribute over
$+$. The same holds for $\rtimes$. Conjecturally, neither tensor is closed.
\prov{MacBeth (computed); closure Open}
\end{Proposition}

\begin{proof}
The shape set of $p+p'$ is $S_p+S_{p'}$, so
$\bigl((p+p')\ltimes q\bigr)[(s,t)]$ has the exponent
$q[t]^{\,S_p+S_{p'}}\cong q[t]^{S_p}\times q[t]^{S_{p'}}$, whereas the
$p$-summand of $(p\ltimes q)+(p'\ltimes q)$ has exponent $q[t]^{S_p}$. These
differ (for a concrete instance the direction cardinality profiles are
$[8,8,32,32,32,32]$ against $[8,8,16,16,16,16]$), so there is no natural
distributive isomorphism. The obstruction is exactly the exponent's dependence
on the \emph{entire} shape set $S_p$, which a coproduct enlarges. A tensor whose
left action fails to preserve coproducts cannot have a right adjoint in that
variable, so $(-)\ltimes q$ is not a left adjoint and $\ltimes$ is not
left-closed; we record the closure of $\rtimes$ as conjecturally negative for
the same reason.
\end{proof}

These are, to our knowledge, the first non-closed monoidal structures in the
$\Cont$ story, in sharp contrast with the census: $\times$ is cartesian closed,
$\dir$ is closed, and $\tri$ has a right coclosure. Because $\ltimes/\rtimes$ are
non-convolutional, this non-closure says nothing about whether the
\emph{convolutional} tensors are closed; the two questions are disjoint.

\subsection{Related work}
\label{sec:dia-related}

We are careful about what is new here, because the theme is well populated. What
this section contributes is an \emph{identification} --- DJN's uninterpreted
$\ltimes$ is de Paiva's Dialectica tensor, extended off the homogeneous slice ---
together with the closed forms of Lemma~\ref{lem:nfold} that expose the
directedness of $\rtimes$, and the taxonomy placement of
§\ref{sec:dia-consequences}. It is emphatically \emph{not} the first connection
between Dialectica and polynomial functors; that relationship is the animating
theme of DJN itself.

The following distinctions matter.

\emph{Lucatelli~Nunes--Vákár \cite{lnv2405}} give a Dialectica-formula monoidal
\emph{closed} structure on Grothendieck constructions, which instantiates to
containers. But their tensor is \emph{an untwisted, product-like tensor lifted
from base and fibre}, and the Dialectica
twist lives in their internal hom $\multimap$, not in the tensor. They neither
produce $\ltimes$ nor address DJN's open question; their construction is
orthogonal to the identification made here.

\emph{Capucci--Gavranović--Malik--Rios--Weinberger \cite{capucci2024}} unify
lenses, optics and Dialectica fibrationally --- at the level of the
\emph{category}, as a common fibrational pattern, not as a monoidal tensor on
$\Cont$. Again this is a different object from $\ltimes$.

Finally, one must pre-empt the objection that DJN's own abstract already
``extends dialectica categories''. It does --- but their extending tensor (their
\S3) is the \emph{non-twisting} parallel product, with direction $A\times B$;
that is the Gödel conjunction, not the multiplicative one. The genuinely
twisting multiplicative tensor, with the mutual function spaces $X^{V}\times
Y^{U}$, is $\ltimes$, and $\ltimes$ is exactly what DJN left uninterpreted. It is
that gap this section fills, framed as an answer to DJN's own stated open
question.

\begin{Remark}[Grade]
\label{rem:dia-grade}
The identification of Proposition~\ref{prop:ltimes-is-dial} is a matching of
definitions and is \emph{proved} at the level of the formulas, not conjectural.
Its \emph{novelty} rests on the observation that DJN, who know
$\Dial(\Set)\simeq\Hmg(2)$ intimately, nonetheless posed the interpretation of
$\ltimes/\rtimes$ as open; a full novelty check against the recent
Dialectica--polynomial literature (de~Paiva, Trotta, Spivak, Hedges, 2023--2026)
is owed before publication. We grade the mathematics \emph{computed/an
identification} and flag the novelty as pending that check.
\end{Remark}

%% ---------------------------------------------------------------------------
\begin{thebibliography}{9}

\bibitem{depaiva89}
V.~de~Paiva,
\emph{The Dialectica Categories},
in \emph{Categories in Computer Science and Logic} (Boulder, 1987),
Contemp.\ Math.\ \textbf{92}, Amer.\ Math.\ Soc., 1989, pp.~47--62;
see also \emph{The Dialectica Categories}, Ph.D.\ thesis / Tech.\ Report
213, Univ.\ of Cambridge Computer Laboratory, 1991.

\bibitem{djn2305}
G.~Dorta, P.~Jarvis, and N.~Niu,
\emph{Monoidal Structures on Generalized Polynomial Categories},
arXiv:2305.05655; EPTCS \textbf{397} (2023).

\bibitem{lnv2405}
F.~Lucatelli~Nunes and M.~Vákár,
\emph{Monoidal closure of Grothendieck constructions via $\Sigma$-tractable
monoidal structures and Dialectica formulas},
arXiv:2405.07724 (2024); to appear in \emph{Theory and Applications of
Categories}.

\bibitem{capucci2024}
M.~Capucci, B.~Gavranović, A.~Malik, F.~Rios, and J.~Weinberger,
\emph{On a fibrational construction for optics, lenses, and Dialectica
categories},
arXiv:2403.16388 (2024); ENTICS Vol.\ 4 (Proc.\ MFPS XL, 2024).

\bibitem{niuspivak}
N.~Niu and D.~I.~Spivak,
\emph{Polynomial Functors: A Mathematical Theory of Interaction},
arXiv:2312.00990 (2023).

\end{thebibliography}

\end{document}
